\chapter{Introduction: Classical error correction and foundations of quantum error correction}
\label{chap:introduction}

The field of quantum computing is evolving rapidly. It has transitioned from a theoretical curiosity to being generally accepted as a frontier technology. The fundamental premise of this discipline is that information processing based on the laws of quantum mechanics can solve problems that remain computationally out of reach for classical computing \cite{nielsen2010quantum}. This paradigm began in 1982, when Richard Feynman pointed out that classical computers are not well equipped to simulate quantum systems due to the exponential growth of space required to represent a quantum state. He proposed that a computer based on quantum mechanical principles would be required to model nature effectively \cite{feynman1982simulating}.

\subsubsection*{The Quantum Advantage and the Noise Barrier}

The true potential of this technology was demonstrated in 1994 by Peter Shor, who developed a quantum algorithm for integer factorization in polynomial time \cite{shor1994algorithms}. Shor's work proved that quantum computing could break widely used cryptographic schemes, such as RSA. That is a theoretical proof of what has been called "Quantum Advantage". However, this assumes the existence of quantum hardware that has not been realized yet.

The realization of such a computer is difficult because of the fragility of quantum states. Qubits are vulnerable to noise from environmental interactions and to inaccuracies in the operations we make with them. When this noise affects the qubits, the information they store is lost.

We refer to this reality as the era of Noisy Intermediate-Scale Quantum (NISQ) technology. In this era, error rates in quantum gates are several orders of magnitude higher than their classical counterparts. Without a robust way to protect quantum information, even the most advanced algorithms fail before reaching a meaningful result.

\subsubsection*{The Evolution of Quantum Error Correction}

As we will see later in this chapter, unlike classical bits, quantum information cannot be protected by just copying it, due to the no-cloning theorem. To overcome this, the field of Quantum Error Correction (QEC) was pioneered by Peter Shor \cite{shor1995scheme} and Andrew Steane \cite{steane1996error}. They demonstrated that a single "logical" qubit can be represented with multiple "physical" qubits. More formally, the logical qubit is encoded in a larger Hilbert space of multiple physical qubits, allowing for the detection and correction of errors without destroying the underlying quantum state.

\notatfm{The Threshold Theorem \cite{aharonov1997fault} states that arbitrarily long quantum computation is only possible if the physical error rate is strictly below a certain critical value, known as the threshold.}
As we will see in the following chapters, while Shor's code proves that QEC can be theoretically implemented, it is unable to effectively mitigate error, as the added complexity of the code causes more errors than it can prevent. This occurs when the physical error rate is above the fault-tolerance threshold, a regime known as "operating above the threshold". The fault-tolerance threshold depends on the specific QEC protocol, including the mechanism to encode the information and detect errors (known as \textbf{syndrome extraction}) and the algorithm used to infer the corrections (known as \textbf{the decoder}). On the other hand, the physical error rate depends on the specific hardware used.

A significant milestone in QEC was the introduction of topological codes by Alexei Kitaev \cite{kitaev2003fault}, which will be discussed in Chapter \ref{chap:state_of_the_art}. Kitaev's Toric Code and subsequent Surface Codes presented a new paradigm of local, lattice-based architectures where errors are detected by measuring local parities (stabilizers). These types of codes became leading candidates for scalable hardware due to their high error thresholds and two-dimensional physical layouts, which map well to the most common realizations of quantum computers.

\subsubsection*{Project Scope}

This Master’s Thesis focuses on the simulation and comparative analysis of decoders for the 2D triangular color code. Introduced by Bombín and Martin-Delgado \cite{bombin_martindelgado_2006}, color codes represent an evolution over surface codes that offers significant advantages, particularly regarding the transversal implementation of logical gates, which simplifies fault-tolerant operations.

The scope of this work is defined by three specific constraints regarding the lattice geometry, the decoding algorithms, and the noise model. The foundations of these concepts are detailed in the following sections and chapters, starting with an overview of classical error correction principles. Below, we outline the specific constraints.

\notatfm{As this introduction establishes the boundaries of the research, it references advanced concepts (like topological geometries or noise models) that will be formally defined in subsequent chapters.} 
First, as we will formally define in Chapter \ref{chap:state_of_the_art}, color codes can be implemented through several geometries (4.8.8, 6.6.6 and 4.6.12). In this work, we will specifically adopt a \textbf{4.8.8 lattice geometry}. This choice is motivated by its lower qubit requirements (fewer qubits are needed per distance) and its properties as a doubly-even code, which enables transversal implementation of the full Clifford group, including $H$, $S^\dagger$ and $\text{CNOT}$ \cite{bombin_martindelgado_2006}.

Second, our objective is to evaluate the performance of various decoding strategies, from classical \textbf{Projection decoders} to state-of-the-art \textbf{Möbius} and \textbf{Concatenated MWPM} approaches. Using advanced simulation and sampling tools such as \texttt{stim} \cite{gidney2021stim} and \texttt{sinter}, in Chapter \ref{chap:comparative_analysis} we will quantify their efficiency through standard metrics: physical error thresholds ($p_{th}$), logical failure rate scaling ($p_L$) as a function of physical noise, theoretical sub-threshold suppression limits, and the trade-off between decoding accuracy and computational time.

Finally, as we will discuss in Chapter \ref{chap:modeling_quantum_error}, various sources of noise affect real quantum hardware. While realistic hardware-aware simulations of QEC codes require accounting for all of them—including imperfect gates and measurement errors—this study focuses strictly on \textbf{code-capacity noise}. This foundational model isolates the decoherence of data qubits over time by assuming perfect syndrome extraction. Adopting this approach allows us to evaluate the strengths of the 4.8.8 lattice and the algorithmic efficiency of the decoders, without the confounding variables introduced by circuit-level operations.

\section{Classical error correction}
\label{sec:classical_error}

Before discussing how to protect quantum information, it is essential to understand the basic mechanisms of classical error correction. While classical bits are physically robust, they are still susceptible to occasional errors (e.g. caused by hardware degradation or electromagnetic interference). Since classical systems are naively discrete (information is binary), a classical error is strictly discrete: a bit flips from $0$ to $1$, or from $1$ to $0$.

To detect and correct these bit-flips, classical computing relies on the concept of \textbf{redundancy}. By adding extra bits to the original message, we can deduce if an error has occurred and, in many cases, revert it. Two of the most foundational techniques used to achieve this are repetition codes and parity checks. These classical concepts serve as the baseline for building analogous quantum error correction solutions.

\subsubsection*{Repetition Codes and Parity Checks}

The simplest approach to redundancy is the \textbf{repetition code}. If we want to send a single logical bit, we copy it multiple times. For example, in a 3-bit repetition code, a logical $0$ is encoded as $000$, and a logical $1$ as $111$. 

\notatfm{For the 3-bit repetition code to be useful, the probability of a single bit flipping, $p$, must be low enough that the chance of two or three bits flipping simultaneously ($3p^2(1-p) + p^3$) is negligible.}
If a noise event flips one of the physical bits (e.g., the state becomes $010$), the receiver can apply a \textbf{majority vote} to determine the original message. Since there are more zeros than ones, the system concludes the intended logical state was $0$, effectively correcting the error.

% Figure: 3-bit Repetition Code
\begin{figure}[hbt!]
\centering
\resizebox{\textwidth}{!}{
    \begin{tikzpicture}[
        box/.style={draw, rectangle, minimum height=1cm, minimum width=2cm, align=center, font=\sffamily\small},
        arrow/.style={-{Stealth}, thick}
    ]
        % Nodes
        \node (input) {Logical $0$};
        \node[box, right=1cm of input] (encode) {Encoding\\$0 \rightarrow 000$};
        \node[box, right=1cm of encode] (noise) {Noisy Channel\\(Bit-flip: $010$)};
        \node[box, right=1cm of noise] (decode) {Decoding\\(Majority Vote)};
        \node[right=1cm of decode] (output) {Logical $0$};
    
        % Arrows
        \draw[arrow] (input) -- (encode);
        \draw[arrow] (encode) -- node[above] {$000$} (noise);
        \draw[arrow] (noise) -- node[above] {$010$} (decode);
        \draw[arrow] (decode) -- (output);
    \end{tikzpicture}   
}
\caption{Mechanism of a classical 3-bit repetition code correcting a single bit-flip error.}
\label{fig:classical_repetition}
\end{figure}

\notatfm{This overlapping block parity concept is the foundation of classical Hamming codes, which can mathematically pinpoint exactly which bit flipped without needing massive repetition.}
A more efficient method is the use of \textbf{parity checks}. Instead of copying the entire message, we calculate a parity bit based on the sum of the data bits. For instance, we can add a bit to ensure that the total number of $1$s in a string is always an even number. If a single bit flips during transmission, the overall parity becomes odd, immediately alerting the system to the presence of an error. By structuring multiple parity checks across different overlapping blocks of data, classical systems can locate and correct errors with minimal overhead.

% Figure: Parity Check Mechanism
\begin{figure}[hbt!]
\centering
    \begin{tikzpicture}[
        bit/.style={draw, rectangle, minimum size=0.8cm, font=\sffamily},
        parity/.style={draw, rectangle, minimum size=0.8cm, fill=gray!20, font=\sffamily\bfseries},
        arrow/.style={-{Stealth}, thick}
    ]
        % Data bits in a row
        \node[bit] (d1) at (0,0) {1};
        \node[bit, right=0cm of d1] (d2) {0};
        \node[bit, right=0cm of d2] (d3) {1};
        
        % Parity bit in the same line
        \node[parity, right=0.8cm of d3] (p) {0};
        
        % Labels
        \node[above=0.1cm of d2, font=\footnotesize] {Data bits (Message)};
        \node[above=0.1cm of p, font=\footnotesize] {Parity bit};
    
        % Brace and dependency line
        \draw [thick, decoration={brace, mirror, raise=0.2cm}, decorate] 
            (d1.south west) -- (d3.south east) 
            node [pos=0.5, anchor=north, yshift=-0.3cm] (logic) {\footnotesize Parity calculation};
            
        \draw[arrow] (logic.east) -- ++(0.4,0) -- ++(0,0.4) -- (p.west);
        
        % Dashed box for the whole codeword
        \draw[dashed, gray] ($(d1.north west)+(-0.2,0.6)$) rectangle ($(p.south east)+(0.2,-1.1)$);
        \node[below=0.7cm of p, font=\scriptsize\color{gray}] {Data block};
    \end{tikzpicture}
\caption{A parity check bit appended to a 3-bit message. The parity bit (0) is calculated to ensure that the total number of '1's in the data block is even.}
\label{fig:parity_check}
\end{figure}

\subsubsection*{The Three Major Challenges of Quantum Error Correction}

It is natural to assume we could simply apply these classical techniques to quantum computers. However, translating repetition and parity concepts to the quantum realm is not straightforward. Quantum mechanics imposes three fundamental constraints that classical computing does not face:

\begin{enumerate}
    \item \textbf{The No-Cloning Theorem:} Classical repetition codes require copying the information. In quantum mechanics, it is fundamentally impossible to create an identical copy of an unknown arbitrary quantum state \cite{wootters1982single}. Therefore, we cannot simply duplicate a qubit to create redundancy.
    \item \textbf{Measurement Destroys Information:} Classical parity checks require reading the bits to see if an error occurred. If we attempt to measure a qubit in a superposition state to check for errors, the superposition collapses into a definitive classical state. Measuring the data to find the error destroys the quantum information we are trying to protect.
    \item \textbf{Continuous Errors:} Classical computers are binary (a bit is either $0$ or $1$) and discrete (there are no intermediary values). This property applies to classical errors (a bit is either right or wrong). Qubits, however, exist in a continuous state space, the complex Hilbert space. This means that there is an infinite number of possible errors that can occur, making it apparently impossible to design a code that accounts for every potential continuous deviation.

\notatfm{A qubit state, $|\psi\rangle$, is a vector in a two-dimensional complex Hilbert space, defined by continuous amplitudes $\alpha|0\rangle + \beta|1\rangle$, with $\alpha$ and $\beta$ being complex numbers. Such state can be represented in the Bloch Sphere. Because of this continuity, an error can be an arbitrarily small rotation.}
\begin{figure}[hbt!]
\centering
\begin{tikzpicture}[line cap=round, line join=round, >=Stealth, scale=2]
    \tikzstyle{axis} = [->, thick, gray!80]
    \tikzstyle{back_axis} = [dotted,  thin, gray!60] 
    \tikzstyle{vector} = [->, very thick, red]
    \tikzstyle{projection} = [dashed, gray!70, thin] 
    \tikzstyle{label_state} = [font=\small, black]
    
    \draw[gray!30] (0,0) circle (1); 
    \draw[gray!70,  dotted, very thin] (0,0) ellipse (1 and 0.25); 
    \draw[gray!60] (1,0) arc (0:-180:1 and 0.25);      

    \draw[axis] (0,0) -- (0,1.2) node[anchor=south, label_state]{$|0\rangle$};
    \draw[thick, gray!80] (0,0) -- (0,-1.1) node[anchor=north, label_state]{$|1\rangle$};

    \draw[axis] (0,0) -- (-0.75,-0.22) node[anchor=east, label_state, xshift=-1pt]{$|+\rangle$};
    \draw[back_axis] (0,0) -- (0.75,0.22) node[anchor=west, gray!40, font=\footnotesize, xshift=1pt]{$|-\rangle$};

    \draw[axis] (0,0) -- (1.2,0) node[anchor=west, label_state]{$|+i\rangle$};
    \draw[back_axis] (0,0) -- (-1.2,0) node[anchor=east, gray!40, font=\footnotesize]{$|-i\rangle$};

    \coordinate (P) at (0.45, 0.75); 
    \coordinate (P_base) at (0.45, -0.05); 

    \draw[projection] (0,0) -- (P_base);
    \draw[projection] (P_base) -- (P);

    \draw[vector] (0,0) -- (P) node[anchor=south west, black] {$|\psi\rangle$};

    \draw[->, thick, blue] (0,0.5) to [bend left=20] (0.25, 0.42);
    \node[blue, font=\footnotesize] at (0.12, 0.6) {$\theta$};

    \draw[->, thick, blue] (-0.25,-0.07) to [bend right=20] (0.28,-0.04);
    \node[blue, font=\footnotesize] at (0.05, -0.2) {$\phi$};

    \node[draw, fill=white, rounded corners, align=left, font=\small] at (1.7, 0.8) {
        $|\psi\rangle = \alpha|0\rangle + \beta|1\rangle$ \\
        $\alpha = \cos(\frac{\theta}{2})$ \\
        $\beta = e^{i\phi}\sin(\frac{\theta}{2})$
    };
\end{tikzpicture}
\caption{Geometric representation of a qubit on the Bloch sphere. The parameters $\theta$ and $\phi$ define the continuous orientation of the state vector $|\psi\rangle$ relative to the computational basis states $|0\rangle$ and $|1\rangle$.}
\label{fig:bloch_sphere}
\end{figure}
\end{enumerate}

The solution to these three fundamental barriers was formulated by Peter Shor, leading to the creation of the first true quantum error correction code, which we explore in the next section.

\section{Shor codes}
\label{sec:shor_codes}

The transition from classical to quantum error correction was made possible by Peter Shor's 1995 proposal \cite{shor1995scheme}. His approach demonstrated that the three barriers of quantum mechanics mentioned above could be bypassed using entanglement, parity-based measurements and concatenation of QEC codes.

\subsection*{The 3-qubit Code: Solving Cloning and Measurement}

\notatfm{The 3-qubit codes exposed here, which are simple repetition codes, are the main components of the 9-qubit Shor code. It is important to note that in Shor’s original paper \cite{shor1995scheme} these 3-qubit codes were not introduced as independent, standalone protocols.}
The first step is understanding that while we cannot copy an unknown state $|\psi\rangle$, we can \textbf{distribute} its information across an entangled system. In the 3-qubit bit-flip code, the encoding does not create copies; it creates a multi-particle state where the logical information is stored in a correlation of these particles' states.
$$|0\rangle_L = |000\rangle, \quad |1\rangle_L = |111\rangle$$

% 3-qubit Bit-flip Encoding Circuit
\begin{figure}[hbt!]
\centering
\begin{quantikz}
    \lstick{$\alpha|0\rangle + \beta|1\rangle$} & \ctrl{1} & \ctrl{2} & \qw & \rstick[3]{$\alpha|000\rangle + \beta|111\rangle$} \\
    \lstick{$|0\rangle$} & \targ{} & \qw & \qw & \\
    \lstick{$|0\rangle$} & \qw & \targ{} & \qw & 
\end{quantikz}
\caption{Encoding circuit for the bit-flip code. Entanglement distributes the logical information without violating the no-cloning theorem.}
\label{fig:bit_flip_encoding}
\end{figure}

To solve the measurement problem, Shor introduced the concept of \textbf{syndrome extraction}. Instead of measuring each qubit (which would destroy the superposition $\alpha|000\rangle + \beta|111\rangle$), we compare pairs of these qubits to determine whether they are equal or different without gaining any information about their actual state.

\begin{figure}[H]
    \centering
    \begin{quantikz}
        \lstick{Data $q_0$} & \ctrl{3} & \qw      & \qw      & \qw \\
        \lstick{Data $q_1$} & \ctrl{2} & \ctrl{3} & \qw      & \qw \\
        \lstick{Data $q_2$} & \qw      & \ctrl{2} & \qw      & \qw \\
        \lstick{Ancilla $s_A$} & \targ{}  & \qw      & \meter{} & \qw \rstick{bit $a$} \\
        \lstick{Ancilla $s_B$} & \qw      & \targ{}  & \meter{} & \qw \rstick{bit $b$}
    \end{quantikz}
    \caption{Syndrome measurement: The ancilla qubits are activated when the two qubits that they monitor are different. Depending on which of the two ancilla qubits are activated, we can determine which of the qubits has flipped.}
    \label{fig:3-qubit-syndrome}
\end{figure}

As shown in Figure \ref{fig:3-qubit-syndrome}, measuring the ancilla qubits does not provide information about the state of each individual qubit, but it provides information about two qubits being different. The following table shows how the measuring of these ancilla qubits maps to specific identified errors.

\begin{table}[H]
    \centering
    \small
    \renewcommand{\arraystretch}{1.2} 
    \begin{tabular}{@{}ccccl@{}}
        \hline
        \textbf{Syndrome ($s_A$)} & \textbf{Syndrome ($s_B$)} & \textbf{Error Location} & \textbf{Correction} \\ \hline
        0 & 0 & None & No action \\
        1 & 0 & $q_0$ & Apply $X_0$ \\
        0 & 1 & $q_2$ & Apply $X_2$ \\
        1 & 1 & $q_1$ & Apply $X_1$ \\ \hline
    \end{tabular}
    \label{tab:3-qubit-syndrome_table}
\end{table}

\subsection*{A Code to Detect Phase-flip Errors}
The bit-flip code addresses $X$ errors. The other fundamental quantum error is the phase-flip ($Z$), which changes the relative phase (the sign) in a superposition: $Z(\alpha|0\rangle + \beta|1\rangle) = \alpha|0\rangle - \beta|1\rangle$. This is strictly a quantum phenomenon and has no classical counterpart.

\notatfm{Formally, the Hadamard gate maps the standard computational basis $\{|0\rangle, |1\rangle\}$ to the superposition basis $\{|+\rangle, |-\rangle\}$. In terms of operators, $HZH = X$.}
To correct a phase-flip, we can leverage a key property of quantum gates: a Hadamard gate ($H$) transforms a $Z$ (phase-flip) error into an $X$ error (bit-flip). Using the $H$ gate, we can reuse the logic of the bit-flip code, where an $X$-parity check (i.e. comparing pairs of qubits) in the logical space effectively functions as a $Z$-check in the physical space.

The logical states are thus defined as:
\[ |0\rangle_L = |+++\rangle, \quad |1\rangle_L = |---\rangle \]

The encoding circuit entangles the qubits using the CNOT structure of the bit-flip code and then applies a global Hadamard transformation to switch the entire block into the relevant basis.

% Circuit: 3-qubit Phase-flip Encoding using quantikz2
\begin{figure}[H]
\centering
\begin{quantikz}
    \lstick{$\alpha|0\rangle + \beta|1\rangle$} & \ctrl{1} & \ctrl{2} & \gate{H} & \qw & \rstick[3]{$\alpha|+++\rangle + \beta|---\rangle$} \\
    \lstick{$|0\rangle$} & \targ{} & \qw & \gate{H} & \qw & \\
    \lstick{$|0\rangle$} & \qw & \targ{} & \gate{H} & \qw &
\end{quantikz}
\caption{Encoding circuit for the 3-qubit phase-flip code. The information is first distributed via entanglement, and then Hadamard gates transform the final state into the $\{|+\rangle, |-\rangle\}$ basis.}
\label{fig:phase_flip_encoding}
\end{figure}

Applying the same measuring technique to these three qubits in pairs, allows determining the syndrome, which in this case will point to a $Z$ error. The Figure \ref{fig:3-qubit-phase-syndrome} shows the corresponding circuit.

\begin{figure}[H]
    \centering
    \begin{quantikz}
        \lstick{Data $q_0$} & \qw & \targ{} & \qw & \qw & \qw & \qw & \qw \\
        \lstick{Data $q_1$} & \qw & \qw & \targ{} & \targ{} & \qw & \qw & \qw \\
        \lstick{Data $q_2$} & \qw & \qw & \qw & \qw & \targ{} & \qw & \qw \\
        \lstick{Ancilla $s_A$ $\ket{0}$} & \gate{H} & \ctrl{-3} & \ctrl{-2} & \qw & \qw & \gate{H} & \meter{} \rstick{bit $a$} \\
        \lstick{Ancilla $s_B$ $\ket{0}$} & \gate{H} & \qw & \qw & \ctrl{-3} & \ctrl{-2} & \gate{H} & \meter{} \rstick{bit $b$}
    \end{quantikz}
    \caption{Phase-flip syndrome measurement: The ancilla qubits are initialized in the $\ket{+}$ state using Hadamard gates. To perform the measurement in the $X$ basis, CNOT gates are applied in reverse (ancilla as control, data as target) to exploit phase kickback. A final Hadamard gate returns the ancillas to the computational basis before measurement.}
    \label{fig:3-qubit-phase-syndrome}
\end{figure}

The following table shows how each of the possible syndromes points to one specific error correction: 

\begin{table}[H]
    \centering
    \small
    \renewcommand{\arraystretch}{1.2} 
    \begin{tabular}{@{}ccccl@{}}
        \hline
        \textbf{$X_1 X_2$} & \textbf{$X_2 X_3$} & \textbf{Error Location} & \textbf{Correction} \\
        0 & 0 & None & No correction ($I$) \\
        1 & 0 & Qubit 1 & Apply $Z_1$ \\
        1 & 1 & Qubit 2 & Apply $Z_2$ \\
        0 & 1 & Qubit 3 & Apply $Z_3$ \\ \hline
    \end{tabular}
    \label{tab:3-qubit-z-syndrome}
\end{table}

\subsection*{The 9-qubit Shor Code: Correcting Arbitrary Errors}
The 3-qubit code provides a solution to the first two challenges, but it only corrects one type of error (either bit-flip or phase-flip). To be able to correct flips in any of the three axis, Shor's proposed the \textbf{concatenation} of a bit-flip code and a phase-flip code. By nesting them, we create a 9-qubit block that protects against both bit-flips ($X$) and phase-flips ($Z$) simultaneously. 

The logical basis states are constructed by taking three clusters of qubits. Each cluster is in a superposition ($|000\rangle \pm |111\rangle$) to handle phase errors, while each individual qubit within the cluster is tripled to handle bit-flip errors:

$$|0\rangle_L = \frac{(|000\rangle + |111\rangle)(|000\rangle + |111\rangle)(|000\rangle + |111\rangle)}{\sqrt{8}}$$
$$|1\rangle_L = \frac{(|000\rangle - |111\rangle)(|000\rangle - |111\rangle)(|000\rangle - |111\rangle)}{\sqrt{8}}$$

\begin{figure}[hbt!]
    \centering
    \begin{quantikz}
        % --- Fase 1: Codificación de Phase-flip (Bloques) ---
        \lstick{$|\psi\rangle$} & \ctrl{3} & \ctrl{6} & \gate{H} & \ctrl{1} & \ctrl{2} & \qw & \rstick[9]{$\alpha|0\rangle_L + \beta|1\rangle_L$} \\
        \lstick{$|0\rangle$} & \qw & \qw & \qw & \targ{} & \qw & \qw & \\
        \lstick{$|0\rangle$} & \qw & \qw & \qw & \qw & \targ{} & \qw & \\
        \lstick{$|0\rangle$} & \targ{} & \qw & \gate{H} & \ctrl{1} & \ctrl{2} & \qw & \\
        \lstick{$|0\rangle$} & \qw & \qw & \qw & \targ{} & \qw & \qw & \\
        \lstick{$|0\rangle$} & \qw & \qw & \qw & \qw & \targ{} & \qw & \\
        \lstick{$|0\rangle$} & \qw & \targ{} & \gate{H} & \ctrl{1} & \ctrl{2} & \qw & \\
        \lstick{$|0\rangle$} & \qw & \qw & \qw & \targ{} & \qw & \qw & \\
        \lstick{$|0\rangle$} & \qw & \qw & \qw & \qw & \targ{} & \qw & 
    \end{quantikz}
    \caption{Full encoding circuit for the 9-qubit Shor code. The initial CNOT and Hadamard gates encode the logical state into a 3-qubit phase-flip code. Subsequently, each of these three qubits is encoded into a 3-qubit bit-flip code, resulting in the final 9-qubit state.}
    \label{fig:shor_9qubit_circuit}
\end{figure}

\notatfm{Z-type detectors are implemented with CNOT gates that have the ancilla as target and the data qubits as controls. X-type detectors are implemented using phase kickback: CNOT gates that have the ancilla as control and the data qubits as target.}
To extract the syndrome of the 9-qubit Shor code, a total of 8 ancilla qubits are needed, 6 of them monitoring pairs of qubits to detect bit flips inside the clusters (Z-type detectors), and 2 of them monitoring full clusters to detect phase flips between two clusters (X-type detectors). The Table \ref{tab:9-shor-ancilla} shows the ancilla qubits and the data qubits they monitor.

\begin{table}[hbt!]
    \centering
    \renewcommand{\arraystretch}{1.2}
    \begin{tabularx}{\textwidth}{@{}cccX@{}}
    \toprule
    \textbf{Ancilla} & \textbf{Detector type} & \textbf{Targeted qubits} & \textbf{Error detected} \\ \midrule
    $a_1$ & $Z$-type & $q_1, q_2$ & Bit-flip ($X$) in Block 1 \\
    $a_2$ & $Z$-type & $q_2, q_3$ & Bit-flip ($X$) in Block 1 \\
    $a_3$ & $Z$-type & $q_4, q_5$ & Bit-flip ($X$) in Block 2 \\
    $a_4$ & $Z$-type & $q_5, q_6$ & Bit-flip ($X$) in Block 2 \\
    $a_5$ & $Z$-type & $q_7, q_8$ & Bit-flip ($X$) in Block 3 \\
    $a_6$ & $Z$-type & $q_8, q_9$ & Bit-flip ($X$) in Block 3 \\
    $a_7$ & $X$-type & $q_1, q_2, q_3, q_4, q_5, q_6$ & Phase-flip ($Z$) in clusters 1 or 2 \\
    $a_8$ & $X$-type & $q_4, q_5, q_6,  q_7, q_8, q_9$ & Phase-flip ($Z$) in clusters 2 or 3 \\ \bottomrule
    \end{tabularx}
    \caption{Ancilla qubits to extract the syndrome in the 9-qubit Shor code.}
    \label{tab:9-shor-ancilla}
\end{table}

The Table \ref{tab:9-shor-syndrome} shows some examples of possible measured syndromes, the corresponding errors and the necessary recovery operation.

\begin{table}[hbt]
\tablenote{It is important to note that phase-flip correction can be applied to any qubit of the cluster for which we detected the error. This is because the phase is encoded in the relative phase between the $|000\rangle$ and $|111\rangle$ components of the cluster, which changes by changing the phase of any of the three qubits. To correct the Y error, the combination of X and Z can be used.}
    \centering
    \renewcommand{\arraystretch}{1.2}
    \begin{tabularx}{\textwidth}{@{}cccX@{}}
    \toprule
    \textbf{Syndrome ($Z$-type)} & \textbf{Syndrome ($X$-type)} & \textbf{Error Type} & \textbf{Correction} \\
    \textbf{$(a_1 a_2,\space a_3 a_4,\space a_5 a_6)$} & \textbf{$(a_7, a_8)$} &  &  \\
    \midrule
    $(00, 00, 00)$ & $(0, 0)$ & None & $I$ (No error) \\
    \midrule
    $(10, 00, 00)$ & $(0, 0)$ & $X_1$ & Apply $X_1$ \\
    $(11, 00, 00)$ & $(0, 0)$ & $X_2$ & Apply $X_2$ \\
    $(01, 00, 00)$ & $(0, 0)$ & $X_3$ & Apply $X_3$ \\
    \midrule
    $(00, 00, 00)$ & $(1, 0)$ & $Z$ in cluster 1 & Apply $Z_1$ \\
    $(00, 00, 00)$ & $(1, 1)$ & $Z$ in cluster 2 & Apply $Z_4$ \\
    $(00, 00, 00)$ & $(0, 1)$ & $Z$ in cluster 3 & Apply $Z_7$ \\
    \midrule
    $(10, 00, 00)$ & $(1, 0)$ & $Y_1$ & Apply $Y_1$ \\
    \bottomrule
    \end{tabularx}
    \caption{Decoding the syndrome in the 9-qubit Shor code.}
    \label{tab:9-shor-syndrome}
\end{table}

\subsubsection*{The Discretization of Quantum Errors}
We have seen that the 9-qubit Shor code enables correcting both $X$ and $Z$, but how can we correct an infinite number of possible small rotations? 

The solution is surprisingly simple. It relies on the mathematics of quantum measurement. As proven by Knill and Laflamme \cite{knill1997theory}, we do not actually need to correct any arbitrary error. Any arbitrary single-qubit error can be expressed as a linear combination of the Pauli matrices: $I$ (identity, no error), $X$ (bit-flip), $Z$ (phase-flip) and $Y$ (both bit and phase-flip).

By designing a quantum error correction code capable of detecting and correcting independent $X$ and $Z$ errors, such as the one designed by Shor, the act of measuring the error (the syndrome extraction) forces the continuous, arbitrary error to collapse into one of the discrete Pauli operators ($I$, $X$, $Y$ and $Z$). Since $I$ means no error and the Pauli $Y$ operator is simply a combination of $X$ and $Z$ ($Y = iXZ$), \textbf{correcting $X$ and $Z$ is sufficient to correct any arbitrary error}. This phenomenon, known as the discretization of quantum errors, effectively reduces the continuity quantum error problem back to a discrete problem analogous to classical bit-flips.

\subsubsection*{Error Distinguishability}
\label{sec:distinguishability}

The fundamental requirement for a quantum error correction code is that errors that require different recovery operations must result in distinct syndromes. In formal terms, this means that the errors must map the logical state into orthogonal, distinguishable subspaces. If two physical errors produce the same syndrome but result in fundamentally different corrupted states, the decoder will not know which correction to apply, leading to a logical failure.

\notatfm{Any qubit within a cluster suffering a phase-flip error in the 9-qubit Shor code results in the same syndrome. However, an operation on any of the three qubits of the cluster will fix the state of the logical qubit. This is known as degeneracy, and we say that the three phase-flip errors that can affect a cluster are \textbf{degenerated errors}.}
The 9-qubit Shor code works well precisely because it separates bit-flip and phase-flip errors into distinct, non-overlapping parity checks. For example, any bit-flip error in a single qubit will activate either one or two of the first 6 ancillas (see Table \ref{tab:9-shor-syndrome}), which can map uniquely to the specific qubit that flipped. Similarly, any phase-flip error in a single qubit will activate one or two of the last 2 ancillas, which maps to the specific cluster that requires correction.

This ability to uniquely map syndromes to the appropriate recovery operation is what makes a QEC code successful. However, analyzing these subspaces manually becomes unsustainable for larger, more complex quantum codes. This scalability problem motivates the need for a robust, algebraic framework, leading directly to the Stabilizer Formalism that we discuss in the next section.

\subsection{The Stabilizer Formalism}
\label{sec:stabilizer_formalism}

As quantum codes scale, representing the logical state as a vector of $2^n$ amplitudes becomes computationally intractable. Tracking the evolution of a highly entangled state qubit by qubit is unsustainable. The Stabilizer Formalism, introduced by Daniel Gottesman in 1997 \cite{gottesman1997stabilizer}, offers an elegant workaround: instead of describing the quantum state directly, we describe the set of operators that leave the state unchanged. More specifically, we focus on the operators that leave unchanged a valid state (i.e. a state where no error has occurred) but modify an erroneous state.

\subsubsection{Stabilizer Operators}

\notatfm{In linear algebra, an eigenvector $\lvert \psi \rangle$ of an operator $S$ satisfies $S\lvert \psi \rangle = \lambda\lvert \psi \rangle$, where $\lambda$ is a constant called the eigenvalue. In the stabilizer formalism, we strictly look for states where $\lambda = +1$. Because the state remains unchanged, measuring $S$ will yield a deterministic outcome of $+1$, indicating no error. It's important to note that measuring the $+1$  usually results in a $0$ value for the classical syndrome bit.}
The core concept of this framework relies on a simple mathematical condition. We say that an operator $S$ \textbf{stabilizes} a quantum state $\lvert \psi \rangle$ if applying $S$ to the state does not alter it in any way:
\begin{equation}
    S \lvert \psi \rangle = \lvert \psi \rangle
\end{equation}

In practice, a stabilizer operator is simply a combination of Pauli matrices ($X$, $Y$, $Z$) acting on a specific subset of qubits. For example, an operator $Z_1 Z_2$ monitors the phase relationship of the first and second qubits. These operators are what we referred to as "syndrome detectors" or "parity checks" in the previous section when discussing the Shor code.

Note that an error affecting the qubits can also be represented as an operator (e.g. a bit-flip is a $X$ Pauli matrix applied to the flipping qubit).

We can \textbf{measure} a stabilizer operator. Such measurement is implemented as a parity check on the data qubits that are connected to this operator, with the result being reflected in the state of an ancilla qubit. In the prior section, we saw how:
\begin{itemize}
    \item $Z$-type stabilizers, which detect bit-flip errors, are measured with a $\text{CNOT}$ gate that has the data qubits as source and the ancilla qubit as target.
    \item $X$-type stabilizers, which detect phase-flip errors, are measured with a $\text{CNOT}$ gate that has the data qubits as target and the ancilla qubit as source, which produces a \textit{phase kickback}.
\end{itemize}

\notatfm{Mathematically, two operators $A$ and $B$ commute if $AB = BA$, and anti-commute if $AB = -BA$. For Pauli matrices, identical types commute (e.g., an $X$ error commutes with an $X$ stabilizer), but different types acting on the same qubit anti-commute (e.g., an $X$ error anti-commutes with a $Z$ stabilizer).}
We must also introduce the concept of \textbf{commutation}. In the mathematics of Pauli matrices, two operators either commute (meaning their order of application does not matter) or \textbf{anti-commute} (meaning swapping their order introduces a minus sign). Because both stabilizers and errors are Pauli operators, we can determine commutation properties for them.

The algebraic property of commutation is the core mechanism of error detection. If a quantum system is stabilized by an operator $S$, we can periodically measure $S$ to check the health of the qubits connected to it:
\begin{itemize}
    \item If no error has occurred, or if an error occurs that \textit{commutes} with $S$, the state remains a valid $+1$ eigenstate. The measurement returns $+1$, telling the decoder that those specific qubits are safe (or not affected by an error detectable by $S$).
    \item If a physical error occurs that \textit{anti-commutes} with $S$, the error flips the sign of the state. The measurement will return $-1$, acting as a localized alarm that an error has corrupted the monitored qubits.
\end{itemize}

We can see that if we can construct a set of operators that can detect all the possible errors affecting our physical qubits, we can determine the operations required to correct them.

\subsubsection{The Stabilizer Group and Generators}
\notatfm{A \textbf{group} is a set equipped with an operation that combines any two elements to form a third element within the same set. The Pauli group is closed under matrix multiplication (the product of  two Pauli operators is always a Pauli operator).

A group whose elements commute between them is known as an \textbf{abelian group}.}
To build these stabilizer operators, we use the $n$-qubit Pauli group, denoted as $\mathcal{P}_n$. This group consists of all possible tensor products of the standard Pauli matrices ($I, X, Y, Z$) applied to $n$ qubits, multiplied by a global phase factor ($\pm 1, \pm i$). 

Under this formalism, a quantum error correction code is defined by its \textbf{Stabilizer Group} ($\mathcal{S}$), which is a commuting subgroup of $\mathcal{P}_n$, that is, a subset of $\mathcal{P}_n$ for which all its members commute between them. 

The reason why this is a strict requirement is that, for a QEC code to be valid, all the operators in the group must share a common set of eigenvectors with eigenvalue $+1$. These eigenvectors are the logical states of the QEC code, and are not affected by any of the stabilizer operators. Only commuting operators can share a common set of eigenvectors.

Additionally, for $\mathcal{S}$ to define a valid code, the operator $-I$ must not be an element of the group. If $-I$ was an element of $\mathcal{S}$, then we would have $-I \lvert \psi \rangle = \lvert \psi \rangle$. This only holds if $\lvert \psi \rangle = 0$, meaning the code would have no valid quantum states.

We've defined the $\mathcal{S}$ as a set of commuting operators. However, we do not need every single operator in $\mathcal{S}$ to define the code. We only need a small, independent set of operators called \textbf{generators}. Just as a three basis vectors can define an entire 3D space, multiplying these generators together produces all the other stabilizers in the group. This set of independent operators represents the set of detectors that we will have to implement in the quantum circuit.

For an $n$-qubit system, choosing $n-k$ independent commuting generators uniquely defines a QEC code capable of encoding $k$ qubits. Such generators create a protected logical subspace with $2^k$ logical states.

\subsubsection{Standard [[n, k, d]] Notation}

\notatfm{Note that the parameter $n$ does not include the number of ancilla qubits required to extract the syndrome, as it depends on the actual implementation.}
In the field of Quantum Error Correction, stabilizer codes are conventionally described using the $[[n, k, d]]$ notation, which provides a summary of the code's cost and capacity:

\begin{itemize}
    \item \textbf{$n$ (Physical Qubits):} The total number of physical qubits used to construct the code.
    \item \textbf{$k$ (Logical Qubits):} The number of logical qubits encoded.
    \item \textbf{$d$ (Code Distance):} The minimum number of physical errors required to cause an undetected logical failure.
\end{itemize}

Under this convention, the Shor code is formally classified as a $[[9, 1, 3]]$ code, meaning it uses nine physical qubits to protect one logical qubit with a distance of three, and requires eight ($n-k$) stabilizers.

The distance directly determines the error-handling capabilities of the code through two fundamental rules:

\begin{enumerate}
    \item \textbf{Detection:} A code can detect any error affecting up to $d-1$ physical qubits.
    \item \textbf{Correction:} A code can correct up to $t$ physical errors:
    \begin{equation}
        t = \left\lfloor \frac{d - 1}{2} \right\rfloor
    \end{equation}
\end{enumerate}

For the Shor code ($d=3$), these rules imply that the code can detect up to 2 errors and correct $t=1$ error. If two errors occur simultaneously, the syndrome would cause a wrong correction. This is a permanent trade-off in QEC: increasing the correction threshold $t$, requires increasing the distance $d$, which requires more physical qubits $n$.

\subsubsection{The Threshold Theorem}
Proven by Aharonov and Ben-Or \cite{aharonov1997fault}, the Threshold Theorem gives us a framework for arbitrarily long quantum computations. 

We have seen so far that through QEC we can effectively lower down the logical error rate (i.e. the error introduced by the overhead QEC in the circuit is lower than the error effectively corrected, so the resulting logical error is below the physical error rate).

We have also seen how a fundamental parameter of a code is the distance, $d$. It determines the ability of the code to correct simultaneous errors and, therefore, it relates to the capacity of the code to reduce the error rate. And we know that the distance of a code is correlated with the code overhead: to increase the distance, we need to add more physical qubits, $n$, to codify the same number of logical qubits, $k$.

However, as we increase the complexity of quantum computations, we need quantum circuits that are wider (more qubits) and deeper (more steps in the calculation). Because the probability of error is constant through the execution, during the execution of an arbitrarily deep circuit, a non-recoverable error will eventually happen. 

The Threshold Theorem combines all this to provide a foundational proof of the usefulness of Quantum Error Correction. It proves that when we increase the distance of a code, hence introducing overhead but improving the code capacity, the logical noise rate that we obtain behaves differently depending on the physical error rate that our hardware has:

\begin{itemize}
    \item If the physical error rate, $p$ is below a certain value, $p_{th}$, which depends on the code itself, increasing the distance reduces the logical noise. That is, the benefit of increasing the code capacity is greater than the drawback of increasing the circuit overhead.
    \item If $p$ is above $p_{th}$, increasing the distance increases the logical noise. That is, the benefit of increasing the code capacity is smaller than the drawback of increasing the circuit overhead.
\end{itemize}

\notatfm{Note that the threshold, $p_{th}$, of a given code is not an isolated property of the code. It depends on the actual implementation of the code in a certain hardware. In Chapter \ref{chap:state_of_the_art} we will show how mapping a given code to a given system (a set of qubits with a given interconnectivity) adds overhead that change the Threshold of the result.}The value at which the behaviour change is what we call the Threshold, $p_{th}$, of a QEC code. We can conclude that, if we have hardware whose physical error rate $p$ is below the threshold $p_{th}$ of a certain code, we can arbitrarily increase the distance of such code to reduce the effective noise. This compensates the increase in the chances of a non-recoverable error (when a circuit depth increases) by increasing the distance of the code.

In Chapter \ref{chap:implementation_and_simulation}, we will simulate various QEC codes to visualize this phase transition. By plotting the logical error rate $p_L$ as a function of the physical error rate $p$ for different distances, we will numerically identify the threshold $p_{th}$ as the precise crossing point where scaling the distance inverts its effect on logical noise.